% Transcription — fonds Alexandre Grothendieck, shelfmark 40, batch 2 (pages 21-27).
% Established from the facsimile archives/batches/40/batch-02.pdf, cut from a
% copy of the folder fetched through the project's own relay
% (https://grothendieck-relay.onrender.com/source/40.pdf), not uploaded by the
% user; PDF page k is the archivists' page 20+k, with no cover sheet. Per the
% skill transcribe-grothendieck. The folder is twenty-seven pages; this is the
% second and last of its two batches. Batch 1 (pages 1-20) was being
% transcribed in parallel and was not available, so the notation below is
% fixed from these seven pages alone and may need aligning with batch 1.
%
% Pages skipped: none. Pages summarised: none. All seven pages (21-27) are
% his hand, blue ink, fast drafting; page 27 carries only its first ten lines.
% No pagination of his is visible on these leaves, and no date.
%
% What the batch is about. Principal bundles (« fibrés principaux »,
% « espaces homogènes principaux ») under the infinitesimal group B — the
% kernel of Frobenius on the additive group in characteristic p (alpha_p in
% later notation) — over a ring A, described by derivations. Page 21 opens on
% a cancelled proof (boxed and struck through) that the powers u^i,
% 0 <= i <= p-1, are free over A, applying D^alpha to a relation; then, for
% u^p = a, C' = A[U]/(U^p - a) with the derivation D U = 1 is a principal
% B-bundle, the map C' -> C, U |-> u, is a B-homomorphism and hence an
% isomorphism, with a bracketed general argument (a G-homomorphism E -> E' of
% principal bundles is an isomorphism, via the square G x E -> G x E' over
% E x E -> E' x E') and the question whether A must be noetherian. Page 22:
% « Théorème » H^1(A, B) = A/A^p, the class of a in A corresponding to
% C_a = A[U]/(U^p - a) with D_a U = 1, C_a = C_b iff a = b mod A^p; every
% principal B-space is some C_a; the additive law of H^1 is not proved to be
% that of A/A^p. « Corollaire » (pp. 22-23): a commutative A-algebra C carries
% a principal B-structure iff it has a one-element p-basis over A; such
% structures are classes of u mod A, equivalently derivations D with Du = 1;
% sub-A-hyperalgebras U' of the C-hyperalgebra in Hom_A(C, C) with U' = B
% correspond to the classes au + b (a unit, b in A), i.e. to points of the
% projective space P_A(C/A) = (C/A)/A^*. Pages 24-26, « Autre façon d'obtenir
% ces résultats, et des résultats globaux »: on a preschema X over Z/pZ the
% exact sequence 0 -> B_X -> G_{a,X} -> G_{a,X} -> 0 (Frobenius), the
% hyperalgebra O[T] with Delta T = T (x) 1 + 1 (x) T, Phi(T) = T^p, the kernel
% O[T]/T^p O[T]; surjectivity of Phi « in a sufficient technical sense »
% (with a « ? » in the margin); the long cohomology sequence
% 0 -> H^0(X, B_X) -> A -> A -> H^1(X, B_X) -> H^1(X, O_X) -> H^1(X, O_X),
% whence H^1(X, B_X) as an extension of the Frobenius-kernel of H^1(X, O_X) by
% A/A^p, and, for X complete over an algebraically closed field, H^1(X, B_X)
% the subspace of H^1(X, O_X) killed by Frobenius. « Remarque » (pp. 26-27):
% over A = k[[X]], the degree-p radicial covers A[T]/(T^p - X^n); regular
% iff n = 1; for p = 3, n = 2, A[X^{2/3}] is the subring of k[[X^{1/3}]] of
% series with no Y-term, yet as B-covers it is the « sum » of two regular ones
% (adjoining p-th roots of X and of -X + X^2), whose fibre product over A is
% already non-regular (it has nilpotents).
%
% Notation fixed once for the batch. His B for the group is kept as B on pages
% 21-23 (it is sometimes drawn with a doubled stem, e.g. « B-fibré » on
% p. 22); from page 24 the sheaf is the ornate letter, set \mathcal{B}_X, and
% the « \mathcal{B}-revêtements » of page 26 keep it. The C with a doubled
% stroke of pages 22-23 is read as the same algebra C and set C. The additive
% group is set G_{a,X} as he writes it (plain G). The underlined O of the
% structure sheaf is \underline{O}. Frobenius on the hyperalgebra is Phi
% (\Phi), labelled with a lower-case phi over the arrow of page 24 and kept so;
% the arrows of the cohomology sequence carry a bold F, set F.
%
% Readings to check first: the whole boxed paragraph of page 21 (the
% condition after « pourvu que », the underlined « monomis », the parenthesis
% to the right of the square); « espace » in « espace homogène principal »
% (pp. 21-22); the end of p. 22 (« ne paraît être pas facile ») and
% « réduite à un élément »; the symbol before « C » on p. 23 (read \ill{});
% « sous-schéma » and « section unité » on p. 24; the opening of the
% surjectivity claim on p. 25 and its margin; « annulé » on p. 26 and the
% word written above « revêtements ».
% Apparatus: 63 \ill{}, 31 \uncertain{}.
%
% Pass: Opus 5.5 (claude-opus-5-5), 2026-09-26 — first pass, unchecked against the pages by a human.

\documentclass[11pt,a4paper]{article}
% Le préambule commun à toutes les transcriptions du fonds Grothendieck.
%
% Il tient en une idée : **l'incertitude se compose, elle ne se résout pas.**
% Une transcription qui lisse un mot douteux a détruit la seule chose pour
% laquelle on la faisait — un lecteur doit pouvoir savoir, à chaque endroit,
% ce qui était sur le feuillet et ce qui a été ajouté. Les six macros qui
% suivent sont donc l'appareil critique, et non de la décoration.
%
% Les mêmes commandes sont reconnues par `scripts/render.mjs`, qui produit la
% vue de lecture du site. Une macro ajoutée ici doit l'être aussi là-bas,
% faute de quoi le rendu échouera bruyamment — ce qui est voulu : un rendu
% silencieusement faux est pire qu'un rendu absent.

\ProvidesPackage{grothendieck}[2026/08/08 Transcriptions du fonds Grothendieck]

\usepackage{amsmath,amssymb,amsthm}
\usepackage{mathrsfs}
\usepackage{stmaryrd}
% Les diagrammes commutatifs sont partout dans ce fonds : c'est la langue
% même de la géométrie algébrique de Grothendieck, et une transcription qui
% les rendrait en prose ne transcrirait rien.
\usepackage{tikz-cd}
% Français par défaut — c'est la langue des feuillets — mais l'anglais est
% chargé aussi : la traduction et le résumé sont des documents anglais, et la
% typographie française (espaces avant les deux-points) y serait une faute.
% Ces documents basculent par \selectlanguage{english} en tête de corps.
\usepackage[english,french]{babel}
\usepackage{fontspec}
\usepackage{geometry}
\geometry{a4paper,margin=25mm,marginparwidth=28mm}
\usepackage{marginnote}
\usepackage{xcolor}
% ulem, not soul. `soul`'s \st and \ul are built on a character-by-character
% scan that XeLaTeX's font machinery breaks: the document fails with an
% undefined control sequence rather than degrading. ulem does the same two jobs
% and is Unicode-engine safe. `normalem` stops it hijacking \emph, which the
% transcriptions use for Grothendieck's own emphasis.
\usepackage[normalem]{ulem}
\usepackage{hyperref}
\hypersetup{colorlinks=true,linkcolor=black,urlcolor=black,pdfborder={0 0 0}}

% --- Métadonnées du lot -----------------------------------------------------
% Reprises telles quelles par le script de rendu, qui les affiche en tête de
% la vue de lecture. Elles fixent ce que le fichier prétend couvrir : sans
% elles, une transcription flotte sans point d'attache dans un fonds de seize
% mille feuillets.

\newcommand{\folder}[1]{\gdef\grothFolder{#1}}
\newcommand{\batch}[1]{\gdef\grothBatch{#1}}
\newcommand{\leaves}[2]{\gdef\grothFirst{#1}\gdef\grothLast{#2}}
\newcommand{\foldertitle}[1]{\gdef\grothFoldertitle{#1}}
\gdef\grothFolder{?}\gdef\grothBatch{?}\gdef\grothFirst{?}\gdef\grothLast{?}\gdef\grothFoldertitle{}

% --- Le résumé --------------------------------------------------------------
%
% Il ouvre la lecture modernisée, sous le titre « Résumé », et il est composé
% en petit. Ce n'est pas un ornement : c'est le seul passage du document qui
% ne soit pas des mathématiques, et un lecteur doit pouvoir le reconnaître
% avant de l'avoir lu — puis le sauter s'il connaît déjà le sujet, ou s'y
% arrêter s'il ne le connaît pas. Le filet qui le suit marque la reprise du
% corps.

\newenvironment{resume}
  {\small\setlength{\parskip}{0.45em}}
  {\par\medskip\hrule\medskip}

% --- Mots-clés ---------------------------------------------------------------
%
% En anglais, en clôture du résumé de la lecture modernisée : le vocabulaire
% moderne sous lequel chercher ce que les feuillets construisent. Le manifeste
% du site (`npm run manifest`) les extrait pour étiqueter la cote entière —
% cette ligne est la source unique des tags, il n'y a pas de fichier à côté.
\newcommand{\keywords}[1]{\par\smallskip\noindent
  {\footnotesize\textsc{Keywords} --- \textit{#1}}\par}

% --- L'appareil critique ----------------------------------------------------

% Le numéro de feuillet, en marge. C'est l'ancre qui permet de régler un
% désaccord sur une formule en regardant *un* feuillet plutôt qu'une chemise
% de six cents. Le site s'en sert aussi pour tourner les pages du fac-similé
% quand on fait défiler la transcription.
\newcommand{\leaf}[1]{%
  \marginnote{\footnotesize\color{gray!70}\textsf{#1}}%
  \label{leaf:#1}\ignorespaces}

% Illisible : on ne devine pas. Un mot inventé qui se lit comme les autres est
% le pire résultat possible.
\newcommand{\ill}{\textcolor{red!60!black}{[\,\ldots\,]}}

% Lecture douteuse : le mot est donné, et signalé comme incertain.
\newcommand{\uncertain}[1]{\uline{#1}}

% Ajout de l'éditeur — une lettre manquante, un mot sous-entendu.
\newcommand{\add}[1]{\textcolor{blue!50!black}{[#1]}}

% Biffé par Grothendieck lui-même. Ce qu'il a rayé fait partie du document :
% une rature dit souvent où la pensée a changé de direction.
\newcommand{\struck}[1]{\sout{#1}}

% Note du transcripteur, sur l'état matériel du feuillet.
\newcommand{\note}[1]{\par\smallskip{\small\color{gray!60!black}\textit{[#1]}}\par\smallskip}

% Note marginale de Grothendieck, distincte du corps du texte.
\newcommand{\marginal}[1]{\par\smallskip\noindent\hspace{1em}%
  {\small\color{gray!60!black}#1}\par\smallskip}

% --- Titre ------------------------------------------------------------------

\AtBeginDocument{%
  \begin{center}
    {\large\bfseries Fonds Alexandre Grothendieck --- cote n\textsuperscript{o}~\grothFolder}\\[2pt]
    {\itshape\grothFoldertitle}\\[4pt]
    {\small feuillets \grothFirst--\grothLast\ (lot \grothBatch)}\\[2pt]
    {\footnotesize\color{gray!60!black}Transcription établie d'après le fac-similé numérisé
     par l'Université de Montpellier.\\
     Les lectures incertaines sont soulignées, les illisibles notées [\,\ldots\,],
     les ajouts entre crochets.}
  \end{center}
  \vspace{1em}}

\endinput


\folder{40}
\batch{2}
\pages{21}{27}
\dating{[années 1960-1970]}
\watermark{Édition de démonstration}
\foldertitle{Calculs [affines sur des] fibrés principaux affines : notes manuscrites (s.d.)}

\begin{document}

\page{21}

\struck{$A$ est local.} \struck{Montrons d'abord que les $u^i$
($0 \leq i \leq p-1$) sont libres sur $A$. Soit \ill{} \ill{}
$\sum_0^{p-1} \lambda_i u^i = 0$ une relation non triviale entre les $u^i$
($\lambda_i \in A$), et soit $\lambda_\alpha$ le premier coefficient non nul
($0 \leq \alpha \leq p-1$). Appliquons $D^\alpha$, il vient \ill{}
\uncertain{$\alpha!\,\lambda_\alpha = 0$}, d'où $\lambda_\alpha = 0$,
absurde.}\note{tout ce début de page est encadré et barré de trois longues
diagonales ; « $A$ est local » est en outre biffé à part}

Soit $u^p = a$, et soit $C' = A[U]/(U^p - a)$. \struck{\ill{}}
\struck{\ill{} \ill{} \ill{}} Soit $D'$ \add{une dérivation} dans $C'$ définie par
$DU = 1$, alors muni \uncertain{de} $D$, $C'$ devient un \uncertain{espace}
homogène principal sous $B$ [cf ci-dessous] ! Or l'homomorphisme $C' \to C$
défini par $U \mapsto u$, est donc un $B$-homomorphisme. Il en résulte que
\struck{\ill{}} $C' \to C$ est un isomorphisme\note{« une dérivation » est ajouté
au-dessus de la ligne ; « $C' \to C$ est » est écrit au-dessus d'un mot
biffé}

[de façon générale, si $E, E'$ sont deux \uncertain{espaces} fibrés
principaux sous $G$, et \emph{si} \struck{\ill{}} on a un $G$-homom.
$E \to E'$, alors \struck{\ill{} \ill{} \ill{} dans} c'est un isomorphisme,
\uncertain{pourvu que} $G \times E \to G \times E'$ dans la catégorie
$\mathcal{C}$, \ill{} le \uncertain{diagramme} \emph{\uncertain{monomis}}
\ill{}

\begin{tikzcd}
G \times E \arrow[r] \arrow[d, "\wr"'] & G \times E' \arrow[d, "\wr"] \\
E \times E \arrow[r, "\sim"] & E' \times E'
\end{tikzcd}

(\uncertain{comme} \ill{} \uncertain{flèches} : \uncertain{opération} \ill{}
\ill{} \uncertain{loi}\ill{}). Or \struck{et d'autre part
$E \times E \to E' \times E'$ est un}\note{la mise en page du crochet est
confuse : « pourvu que », « le diagramme » et le mot souligné lu
« monomis » sont enfermés dans un cadre qui renvoie au carré, la parenthèse
de droite est cernée d'un trait, et la dernière ligne est biffée ; l'ordre
de lecture proposé est incertain}
c'est le cas ici en prenant pour $\mathcal{C}$ la catégorie des anneaux finis
sur $A$ --- [pourvu que $A$ soit noethérien ? ---]

\page{22}

On a ainsi prouvé le

\emph{Théorème}. $H^1(A, B)$ s'identifie \struck{à $A$} au groupe additif
$A/A^p$ (où $A^p$ est l'ens. des $x^p$, avec $x \in A$). Si \struck{\ill{}}
$a \in A$, il lui correspond la classe
de la $A$-algèbre $C_a = A[U]/(U^p - a)$, munie de la $A$-dérivation $D_a$
définie par $D_a U = 1$, et $C_a \simeq C_b$ si et seulement si
$a \equiv b \ (A^p)$. Soit $C$ un $B$-\uncertain{espace} homogène principal,
alors il existe dans $C$ un \struck{\ill{}} $u$ tel que $Du = 1$, et
\struck{\ill{} $a \in$ \ill{}} on a $u^p = a \in A$, $C \simeq C_a$ en tant
que \struck{fibré} $B$-fibré ; $u$ est bien déterminé mod addition d'un
élément de $A$.\note{un exposant est biffé sur le $A$ de « $a \in A$ »}

On n'a pas prouvé que la loi additive naturelle de $H^1(A; B)$ correspond à
celle de $A/A^p$, mais \ill{} \ill{} \uncertain{paraît} être \uncertain{pas}
facile.

\emph{Corollaire}. Soit $C$ une $A$-algèbre commutative. Pour
\struck{qu'il \ill{}} qu'il existe sur $C$ une structure de fibré \add{homogène} principal sous
$B$, il faut
et il suffit que $C$ ait une $p$-base sur $A$ \uncertain{réduite} \uncertain{à}
un élément, i.e. un $u \in C$ tel que les $u^i$ ($0 \leq i \leq p-1$)
forment une base de $C$ sur $A$. La donnée d'une telle structure de
$B$-fibré \add{homogène} principal sur $C$\note{« qu'il » est récrit au-dessus du mot biffé ;
« homogène » est ajouté deux fois au-dessus de la ligne, avec un renvoi}

\page{23}

équivaut à la donnée d'une classe de tels $u$ mod addition d'un élément de
$A$, i.e. à la donnée d'un élément de $C/A$ tel que, si $u$ en est un
représentant, $(u^i)$ forme une $p$-base de $C$ sur $A$ ; choisi $u$, la
$B$-structure correspondante est définie par la $A$-dérivation $D$ de $C$
telle que $Du = 1$, et inversement, donnée $D$, $u$ est déterminé (mod $A$)
par la condition $Du = 1$.

Enfin, la donnée d'une \struck{restriction des scalaires}
sous-$A$-hyperalgèbre
$\mathcal{U}'$ de
la $C$-hyperalgèbre \uncertain{formée} \uncertain{des}
$\mathrm{Hom}_A(C, C)$, telle que \ill{}${}_A$ $C$ devienne
\struck{\ill{}} homogène principal sous $\mathcal{U}'$ et que
$\mathcal{U}' \simeq B$ --- équivaut à la donnée d'une
\struck{\ill{} $\in$ $p$-base de $u$} classe de $p$-bases $(u)$ réduites à un élément $u$, \ill{}
\uncertain{multiplicativement} \ill{}, la \ill{} donnée par $u$ \ill{} les
éléments $v = au + b$, avec $a \in A^*$ (unités de $A$) et $b \in A$ ; i.e.
l'une des ces données est en correspondance biunivoque avec les éléments
de l'espace \add{projectif}
$\mathbf{P}_A(C/A) = (C/A)/A^*$ associé au $A$-module $C/A$.\note{« sous-$A$-hyperalgèbre », « classe de » et « projectif » sont
écrits au-dessus de la ligne, les deux premiers sur des mots biffés ; le
$\mathcal{U}'$ qui suit « sous-$A$-hyperalgèbre » est porté dans la marge de
gauche}

\page{24}

Autre façon d'obtenir ces résultats, et des résultats \emph{globaux}.

Soit $X$ un \struck{schéma} préschéma
sur $\mathbb{Z}/p\mathbb{Z}$ ($p$ un nb premier).\note{« préschéma » est
écrit au-dessus du mot biffé} On a sur $X$ une suite
exacte de faisceaux de groupes commutatifs
\[
0 \longrightarrow \mathcal{B}_X \longrightarrow G_{a,X}
\xrightarrow{\ \Phi\ } G_{a,X} \longrightarrow 0
\qquad \Phi = \text{Frobenius}
\]
$G_{a,X}$ correspond à l'hyperalgèbre $\underline{O}[T]$ munie de
\[
\Delta T = T \otimes 1 + 1 \otimes T .
\]
L'homom. de Frobenius sur \struck{cette} l'hyperalgèbre est donné par
\[
\underline{O}[T] \xleftarrow{\ \varphi\ } \underline{O}[T]
\]
\struck{\ill{}}
\[
\Phi(T) = T^p
\]
[on vérifie de suite que c'est compatible avec $\Delta$]. Son « noyau » est
le \uncertain{sous-schéma} de $G_{a,X}$ « image inverse » par $\Phi$ de la
\uncertain{section} unité de $G_{a,X}$, \struck{\ill{}} i.e. correspond au
quotient $\underline{O}[T]/\Phi(\mathrm{Ker}\,\varepsilon)$ où
$\varepsilon$ est l'augmentation dans $\underline{O}[T]$, donc ce noyau est
bien $\underline{O}[T]/T^p\underline{O}[T]$ (car $\mathrm{Ker}\,\varepsilon$
est engendré par $T^p$). Je dis
\ill{} que \struck{$G$} $\Phi \colon G_{a,X} \to G_{a,X}$ est
\emph{surjectif}\note{« engendré par $T^p$ » est lu tel quel ; c'est $T$ qui
engendre $\mathrm{Ker}\,\varepsilon$, et $T^p = \Phi(T)$ qui engendre
$\Phi(\mathrm{Ker}\,\varepsilon)$}

\page{25}

en un sens technique suffisant : tout $X$-\uncertain{morphisme}
$G_{a,X} \to Y$ (où $Y$ est un préschéma sur $X$) \uncertain{qui est} un
$G_{a,X}$-hom. (quand $G_{a,X}$ opère sur le premier facteur grâce à $\Phi$,
et sur $Y$ trivialement) est une \emph{section} de $Y$, i.e. se factorise
en
\[
G_{a,X} \xrightarrow{\ \mathrm{proj}\ } X \longrightarrow Y .
\]
Cela se voit p.ex. sur les \struck{faisceaux} \ill{} d'hyperalgèbres $\Phi$, en montrant que
\uncertain{admettant} \ill{} \ill{} est injectif [du moins si $Y$ est
affine ; \struck{\ill{} \ill{}}, \ill{} \ill{} dans ce
cas].\marginal{?}\note{un gros point d'interrogation dans la marge de
gauche, en regard du crochet ; le mot illisible après « faisceaux » est
écrit au-dessus de ce mot biffé}

On en conclut \ill{} la \struck{\ill{}} suite exacte de
cohomologie, qui s'écrit\marginal{: \uncertain{justifier} ?
\ill{}}\note{dans la marge de gauche, écrit de biais, avec une flèche vers
le texte ; le passage est marqué d'un trait vertical}
\struck{$0 \to A[T]/T^p \to$}
\[
0 \to H^0(X, \mathcal{B}_X) \to A \xrightarrow{\ F\ } A \longrightarrow
H^1(X, \mathcal{B}_X) \to H^1(X, \underline{O}_X) \xrightarrow{\ F\ }
H^1(X, \underline{O}_X)
\]
\note{le premier départ de la suite, biffé, est au-dessus de la ligne
définitive ; la première étiquette $F$ récrit une étiquette biffée}
[où $\Phi$ provient de Frobenius \struck{\ill{}}, est \ill{} \ill{} pour
$A$ : $H^1(X, \underline{O}_X)$] d'où $H^1(X, \mathcal{B}_X)$ comme
extension \ill{} \struck{\ill{}} du sous-groupe
\struck{\ill{} \ill{} $H^1(X, \underline{O}_X)$ \ill{}} par $A/A^p$. Si
p.ex. $X$ est une variété complète sur un corps algébriquement clos, on
obtient que

\page{26}

$H^1(X, \mathcal{B}_X)$ est le \uncertain{sous-espace} vectoriel de
l'espace de dim.\ finie $H^1(X, \underline{O}_X)$, \uncertain{annulé} par le
Frobenius.

\emph{Remarque}. Soit $A = k[[X]]$, où $k$ est un corps algébriquement clos.
Les revêtements \struck{\ill{}} \ill{} \struck{ramifiés} de $A$ obtenus par extraction d'une racine $p$-ième
de $A$ sont \ill{} \uncertain{isomorphes} [\uncertain{à condition de}
changer l'uniformisante, au besoin] à
\[
A[X^{n/p}] \simeq A[T]/(T^p - X^n),
\]
où \ill{}${} \leq n \leq p-1$ \struck{\ill{}} [pour $n = 0$, on
obtiendrait $\simeq A[T]/T^pA[T]$]. Si $n = 1$, c'est un anneau local
régulier. Sinon, il ne l'est pas, comme on voit facilement. Par exemple, si
$p = 3$, \struck{\ill{} $A$ \ill{}} et si on prend $n = 2$, on obtient
\[
A[T]/(T^3 - X^2) = A[X^{2/3}],
\]
qui est isomorphe au sous-anneau de \struck{\ill{}}
$k[[Y]] = k[[X^{1/3}]]$ formé des séries formelles dont le coefficient de
$Y$ est nul. Pourtant, \ill{} \uncertain{point} \ill{} \uncertain{vue} des
$\mathcal{B}$-\struck{fibrés} revêtements,\note{« revêtements » est écrit au-dessus du mot biffé ; plus
haut, le mot illisible après « Les revêtements » est écrit au-dessus d'un
mot biffé}

\page{27}

le revêtement envisagé est la somme de deux revêtements réguliers,
correspondants resp.\ à l'adjonction d'une racine $p$-ième de $X$, et de
$-X + X^2$ ; mais le produit fibré de ces deux anneaux
$C_1, C_2$ sur $A$ est déjà un anneau non régulier [car isomorphe à
$C_1[T]/(T^p)$, qui a des éléments nilpotents]\note{« $C_1, C_2$ » est
ajouté au-dessus de la ligne}

\end{document}
